\documentclass[../DoAn.tex]{subfiles}
\begin{document}

% Chương này có độ dài tối thiểu 5 trang, tối đa không giới hạn.\footnote{Trong trường hợp phần này dưới 5 trang thì sinh viên nên gộp vào phần kết luận, không tách ra một chương riêng rẽ nữa.} Sinh viên cần trình bày tất cả những nội dung đóng góp mà mình thấy tâm đắc nhất trong suốt quá trình làm ĐATN. Đó có thể là một loạt các vấn đề khó khăn mà sinh viên đã từng bước giải quyết được, là giải thuật cho một bài toán cụ thể, là giải pháp tổng quát cho một lớp bài toán, hoặc là mô hình/kiến trúc hữu hiệu nào đó được sinh viên thiết kế.

% Chương này \textbf{là cơ sở quan trọng} để các thầy cô đánh giá sinh viên. Vì vậy, sinh viên cần phát huy tính sáng tạo, khả năng phân tích, phản biện, lập luận, tổng quát hóa vấn đề và tập trung viết cho thật tốt.
% Mỗi giải pháp hoặc đóng góp của sinh viên cần được trình bày trong một mục độc lập bao gồm ba mục con: (i) dẫn dắt/giới thiệu về bài toán/vấn đề, (ii) giải pháp, và (iii) kết quả đạt được (nếu có).

% Sinh viên lưu ý \textbf{không trình bày lặp lại nội dung}. Những nội dung đã trình bày chi tiết trong các chương trước không được trình bày lại trong chương này. Vì vậy, với nội dung hay, mang tính đóng góp/giải pháp, sinh viên chỉ nên tóm lược/mô tả sơ bộ trong các chương trước, đồng thời tạo tham chiếu chéo tới đề mục tương ứng trong Chương 5 này. Chi tiết thông tin về đóng góp/giải pháp được trình bày trong mục đó.

% Ví dụ, trong Chương 4, sinh viên có thiết kế được kiến trúc đáng lưu ý gì đó, là sự kết hợp của các kiến trúc MVC, MVP, SOA, v.v. Khi đó, sinh viên sẽ chỉ mô tả ngắn gọn kiến trúc đó ở Chương 4, rồi thêm các câu có dạng: ``Chi tiết về kiến trúc này sẽ được trình bày trong phần 5.1". 
\section{Tối ưu thời gian phản hồi khi xem danh sách sản phẩm}
\subsection{Đặt vấn đề}
Đây là tính năng rất quan trọng, có thể coi là được người dùng thường xuyên, liên tục với tần suất nhiều nhất trong hệ thống, nên tốc độ phản hồi ảnh hưởng rất lớn đến trải nghiệm người dùng với ứng dụng. Về dữ liệu sản phẩm thì trong hệ thống tôi đã thu thập được từ việc crawl dữ liệu trên mạng lại khá lớn, khoảng 4500 sản phẩm khác nhau. Dữ liệu được lưu trữ dưới dạng document với nhiều trường phức tạp, phải gọi nhiều API khác nhau để lấy được đầy đủ các thông tin về sản phẩm đó như thông tin shop, thông tin brand, thông tin categories. Vì có một số lượng lớn data lớn như vậy và mỗi sản phẩm lại phải call nhiều API nên khiến thời gian chờ người dùng khá lâu, ảnh hưởng không tốt đến trải nghiệm người dùng. 

\subsection{Giải pháp}
Để cải thiện tốc độ phản hồi của hệ thống, em thực hiện đồng thời một số giải pháp sau:

Đầu tiên em không load và tải hết cả 4500 sản phẩm của hệ thống, hoặc vài trăm đến vài nghìn sản phẩm thuộc cùng 1 danh mục hay 1 brand lên ngay, vì làm như vậy trung bình hệ thống mất khoảng 10-15s để phản hồi. Thay vào đó em áp dụng chiến lược phân trang, mỗi lần load chỉ load 20 sản phẩm lên, nếu người dùng hứng thú sẽ click nút loadmore ở cuối trang, hệ thống sẽ load thêm 20 sản phẩm tiếp theo, và cứ như vậy cho tới khi hết sản phẩm thì sẽ thay nút loadmore bằng thông báo đã hết sản phẩm. Để thực thi được điều đó thì em sử dụng cơ chế phân trang dạng con trỏ (cursor-based pagination) của Firestore, cho phép truy vấn chính xác các tài liệu tiếp theo dựa trên document cuối cùng được lưu lại. Việc kết hợp với giới hạn truy vấn chỉ trả về 20 sản phẩm giúp hệ thống chỉ tải đúng 20 sản phẩm mỗi lần, tránh dư thừa dữ liệu và tăng tốc độ phản hồi đáng kể. Bằng cách này, hệ thống vừa đảm bảo hiệu năng, vừa mang lại trải nghiệm người dùng mượt mà hơn, đặc biệt trong các trường hợp người dùng duyệt nhiều sản phẩm trong cùng một danh mục hoặc thương hiệu.

Thứ hai, để tối ưu thời gian phản hồi khi hiển thị thông tin chi tiết của từng sản phẩm, em cũng thực hiện song song nhiều truy vấn không phụ thuộc nhau bằng cách sử dụng kỹ thuật bất đồng bộ với \texttt{Future.wait}. Thay vì truy xuất tuần tự dữ liệu của thương hiệu, cửa hàng và danh mục -- mỗi bước đều phải chờ hoàn tất truy vấn trước đó -- em tiến hành khởi tạo trước các \texttt{Future} tương ứng với mỗi nguồn dữ liệu nếu có, và sau đó dùng \texttt{Future.wait} để thực hiện tất cả truy vấn một cách song song. Cách làm này giúp tiết kiệm đáng kể thời gian chờ đợi tổng thể, đặc biệt trong trường hợp người dùng duyệt nhiều sản phẩm liên tục, nhờ đó tăng hiệu suất xử lý và giảm cảm giác trễ khi hiển thị dữ liệu. 

\begin{figure}[H]
    \centering
    \includegraphics[width=1\linewidth]{Hinhve/C5/future_wait.jpg}
    \caption{So sánh khi sử dụng hàm Future.wait so với khi gọi API tuần tự}
    \label{fig:promise}
\end{figure}
\FloatBarrier
\vspace{-\baselineskip}
Hình \ref{fig:promise} biểu diễn sự so sánh giữa khi sử dụng hàm Future.wait so với khi gọi API tuần tự. Như trong hình ta có thể thấy, khi sử dụng Future.wait thời gian chạy của toàn bộ tất cả các tác vụ bằng thời gian chạy của tác vụ có thời gian lớn nhất. Trong khi đó, nếu chỉ thực hiện các tác vụ gọi API theo kiểu tuần tự thì thời gian là bằng tổng thời gian của tất cả các tác vụ cộng lại. Điều đó cho thấy sự hiệu quả khi áp dụng kỹ thuật lập trình bất đồng bộ sử dụng hàm Future.wait khi cần phải xử lý nhiều tác vụ gọi API cùng một lúc.

Thứ ba là với những dữ liệu nặng nhưng lại ít thay đổi, được gọi nhiều lần
ở nhiều vị trí trong ứng dụng như top category, top shop hay top brand thì em chỉ gọi 1 lần duy nhất, và lưu vào cached, những lần sau sẽ check từ trong cached đã có rồi thì sẽ không truy vấn database nữa mà lấy ngay từ cached, giảm thời gian phản hồi của hệ thống.
\subsection{Kết quả}
\begin{center}
    \includegraphics[width=0.6\linewidth]{Hinhve/C5/toi_uu_phan_hoi.jpg}
    \captionof{figure}{Minh họa thời gian phản hồi của hệ thống sau tối ưu}
    \label{fig:optimize_response}
\end{center}
Hình \ref{fig:optimize_response} cho thấy rất rõ, nhờ việc thực hiện đồng bộ các giải pháp trên mà thời gian phản hồi trung bình đã giảm 10 lần, từ khoảng 10-15s xuống còn từ 1-2s với lần load đầu tiên và các lần load sản phẩm tiếp theo, cải thiện rất nhiều về trải nghiệm người dùng và tạo độ mượt mà ngang hệ thống lớn như shopee, tiktok shop.

\section{Triển khai tính năng "Gợi ý sản phẩm cá nhân hóa theo người dùng"}
\subsection{Đặt vấn đề}

Trong một hệ thống thương mại điện tử hoặc bán hàng trực tuyến, số lượng sản phẩm có thể lên tới hàng nghìn, thậm chí hàng chục nghìn mặt hàng thuộc nhiều danh mục và thương hiệu khác nhau. Điều này khiến người dùng, đặc biệt là những người mới, dễ rơi vào tình trạng “quá tải lựa chọn” (choice overload), khi phải mất nhiều thời gian để tìm được sản phẩm phù hợp với sở thích hoặc nhu cầu thực tế của bản thân. Việc để người dùng tự tìm kiếm, lọc và duyệt qua hàng loạt sản phẩm không chỉ làm giảm trải nghiệm sử dụng, mà còn ảnh hưởng tiêu cực tới tỷ lệ chuyển đổi đơn hàng của hệ thống.

Bên cạnh đó, dữ liệu hành vi của người dùng (như lịch sử mua hàng, lượt xem sản phẩm, lượt thêm vào giỏ hàng, v.v.) thường xuyên thay đổi và gia tăng theo thời gian. Do đó, nếu không có cơ chế cập nhật mô hình gợi ý một cách thường xuyên thì các đề xuất sẽ nhanh chóng trở nên lỗi thời, không còn phù hợp với sở thích hiện tại của người dùng. Việc huấn luyện lại mô hình toàn bộ từ đầu mỗi lần cập nhật cũng là giải pháp tốn kém về thời gian và tài nguyên tính toán.

Chính vì vậy, cần thiết phải xây dựng một tính năng gợi ý sản phẩm cá nhân hóa theo từng người dùng dựa trên mô hình học máy, lại có khả năng học gia tăng theo dữ liệu mới để không phải train lại toàn bộ data từ đầu. 

\subsection{Giải pháp}

Để giải quyết vấn đề trên, em lựa chọn sử dụng mô hình LightFM – một mô hình hybrid kết hợp giữa collaborative filtering và content-based filtering – để khai thác cả hành vi người dùng và đặc điểm sản phẩm, từ đó tạo ra các đề xuất chính xác và có giá trị hơn. Đặc biệt, em thiết kế cơ chế tự động huấn luyện lại mô hình theo lịch định kỳ (dựa vào Cloud Scheduler), sử dụng phương pháp huấn luyện gia tăng (incremental learning) thông qua hàm \texttt{fit\_partial} của LightFM. Điều này giúp hệ thống luôn cập nhật được dữ liệu người dùng mới mà không cần huấn luyện lại toàn bộ mô hình từ đầu, vừa tiết kiệm tài nguyên, vừa đảm bảo khả năng thích ứng linh hoạt theo thời gian.

Giải pháp được triển khai theo một quy trình tổng thể, gồm nhiều bước liên kết chặt chẽ từ frontend Flutter đến backend xử lý mô hình gợi ý. Ở phía ứng dụng Flutter, em xây dựng các hàm ghi log để ghi lại những tương tác thực tế của người dùng như xem sản phẩm, thêm vào giỏ, mua hàng, viết bình luận ... Các tương tác này được lưu vào bảng \ref{tab:userEventsTable} trong Firestore, mỗi loại sự kiện sẽ có một trọng số khác nhau phản ánh mức độ quan trọng của event khác nhau.

Tiếp theo, em sử dụng môi trường xử lý dữ liệu trên Kaggle để viết pipeline train mô hình. Pipeline bắt đầu bằng việc đọc dữ liệu từ Firestore, xây dựng ma trận tương tác giữa người dùng và sản phẩm với các trọng số tương ứng cho từng loại sự kiện. Đồng thời, em trích xuất thêm thông tin từ bảng \ref{tab:usersTable} và bảng \ref{tab:productsTable} để xây dựng các đặc trưng nội dung (user features và item features), giúp tăng độ chính xác của mô hình hybrid.

Để tối ưu hiệu quả huấn luyện, em tích hợp thư viện \texttt{Optuna} nhằm tự động tìm ra tổ hợp siêu tham số tốt nhất cho mô hình LightFM, bao gồm cả số lượng epoch, learning rate, item\_alpha, user\_alpha và loss function. Sau khi xác định được bộ tham số tối ưu, em tiến hành huấn luyện mô hình chính thức và lưu mô hình vào file định dạng \texttt{.pkl}, sau đó tải lên Firebase Storage để sử dụng sau này.

Quy trình huấn luyện lại định kỳ được triển khai thông qua một API sử dụng Flask, được deploy trên Google Cloud Run và thiết lập lịch gọi tự động mỗi 24 giờ nhờ Cloud Scheduler. Mỗi lần gọi, API sẽ truy vấn các tương tác mới phát sinh trong bảng \ref{tab:userEventsTable} kể từ lần huấn luyện cuối cùng (dựa vào trường \texttt{last\_trained}). Sau đó, hệ thống tạo lại một ma trận tương tác mới chỉ chứa các dữ liệu mới này, và tiến hành cập nhật mô hình hiện tại bằng phương pháp huấn luyện gia tăng thông qua hàm \texttt{fit\_partial}. Mô hình sau khi được cập nhật sẽ lại được lưu vào Firebase Storage, đảm bảo hệ thống luôn có phiên bản mới nhất phục vụ cho việc gợi ý.

Bên cạnh đó, em xây dựng thêm một API thứ hai (cũng sử dụng Flask) để ứng dụng Flutter có thể gửi request đến và nhận danh sách sản phẩm được đề xuất dưới dạng JSON. Danh sách này sau đó sẽ được hiển thị trực tiếp trên giao diện người dùng, mang lại trải nghiệm gợi ý sản phẩm cá nhân hóa theo thời gian thực. Cả hai API đều được triển khai bằng Google Cloud Run, giúp hệ thống dễ mở rộng, dễ bảo trì, và tích hợp mượt mà với hạ tầng đám mây hiện tại.

Tổng thể, giải pháp này vừa tận dụng được các kỹ thuật học máy hiện đại để mang lại khả năng gợi ý chính xác, vừa đảm bảo hiệu suất và tính mở rộng nhờ thiết kế một pipeline hiện đại. Quá trình xử lý được tự động hóa từ khâu ghi nhận dữ liệu, huấn luyện mô hình, cập nhật định kỳ cho tới triển khai và tích hợp vào ứng dụng, mang lại một hệ thống gợi ý hoàn chỉnh, thông minh và thích nghi tốt với sự thay đổi hành vi người dùng theo thời gian.

\subsection{Kết quả}

Sau khi triển khai tính năng gợi ý sản phẩm cá nhân hóa theo người dùng, hệ thống đã ghi nhận được một lượng dữ liệu tương tác đáng kể từ phía người dùng thực tế, đồng thời các mô hình học máy cũng cho thấy hiệu quả rõ rệt khi so sánh giữa các phiên bản có và không fine-tune, có và không sử dụng meta-data. Kết quả cụ thể được trình bày như sau:

\subsubsection{Thống kê dữ liệu thu thập được}

\begin{center}
    \includegraphics[width=0.6\linewidth]{Hinhve/C5/so_user.jpg}
    \captionof{figure}{Số lượng người dùng đã ghi nhận tương tác}
    \label{fig:so_user}
\end{center}

\begin{center}
    \includegraphics[width=0.6\linewidth]{Hinhve/C5/so_product.jpg}
    \captionof{figure}{Tổng số lượng sản phẩm trong hệ thống}
    \label{fig:so_product}
\end{center}

\begin{center}
    \includegraphics[width=0.6\linewidth]{Hinhve/C5/tong_so_tuong_tac.jpg}
    \captionof{figure}{Tổng số tương tác được ghi nhận trong bảng \texttt{user\_events}}
    \label{fig:tong_so_tuong_tac}
\end{center}

\begin{center}
    \includegraphics[width=0.6\linewidth]{Hinhve/C5/thong_ke_tuong_tac.jpg}
    \captionof{figure}{Tỷ lệ các loại tương tác: view, click, add to cart, purchase}
    \label{fig:thong_ke_tuong_tac}
\end{center}

Các dữ liệu về người dùng như ở hình \ref{fig:so_user} gồm 64 người dùng và dữ liệu mua sắm, tương tác vơi hệ thống của họ ở hình \ref{fig:thong_ke_tuong_tac} đều là các người dùng và tương tác thật từ những người dùng đã thực tế trải nghiệm ứng dụng của em, đảm bảo tính chính xác, khách quan và gợi ý chính xác cho từng cá nhân.

\subsubsection{Cơ chế tự động huấn luyện định kỳ}

Mô hình được huấn luyện định kỳ mỗi 24 giờ để cập nhật các dữ liệu tương tác mới thông qua cơ chế incremental learning nhờ mô hình có hàm fit\_partial cho phép train tiếp liên tục với data bổ sung, đây cũng là 1 điểm mạnh của mô hình này so với các mô hình học máy khác. Việc này được thực hiện nhờ sự kết hợp giữa Flask API, Google Cloud Run và Cloud Scheduler lên lịch tự động.

\begin{center}
    \includegraphics[width=0.6\linewidth]{Hinhve/C5/job_daily.jpg}
    \captionof{figure}{Lên lịch job huấn luyện định kỳ bằng Cloud Scheduler}
    \label{fig:scheduler}
\end{center}

\subsubsection{Kết quả tìm kiếm tham số tối ưu}

Để nâng cao độ chính xác của mô hình, em sử dụng thư viện Optuna để thử nghiệm và chọn ra bộ tham số tối ưu nhất.

\begin{center}
    \includegraphics[width=0.6\linewidth]{Hinhve/C5/best_param.jpg}
    \captionof{figure}{Bộ tham số tối ưu thu được từ Optuna}
    \label{fig:best_param}
\end{center}

\subsubsection{So sánh hiệu quả mô hình có và không fine-tune}

Dưới đây là kết quả so sánh giữa hai phiên bản mô hình: một phiên bản chỉ train ban đầu, và một phiên bản được train bằng bộ tham số tối ưu mà optuna tìm ra, kết quả bộ tham số tối ưu cho ra kết quả vượt trội hơn ở tất cả các chỉ số đánh giá.

\begin{center}
    \includegraphics[width=0.6\linewidth]{Hinhve/C5/so_sanh_2_model.jpg}
    \captionof{figure}{So sánh kết quả mô hình với và không sử dụng fine-tune}
    \label{fig:so_sanh_finetune}
\end{center}

\subsubsection{Hiệu quả của mô hình khi không sử dụng metadata}

Ngoài ra, để kiểm chứng mức độ đóng góp của metadata (đặc trưng sản phẩm và người dùng), em cũng thực hiện huấn luyện mô hình LightFM chỉ sử dụng ma trận tương tác. Kết quả bất ngờ cho thấy mô hình không sử dụng metadata lại cho kết quả tốt hơn về mọi mặt:

\begin{center}
    \includegraphics[width=0.6\linewidth]{Hinhve/C5/mo_hinh_ko_meta.jpg}
    \captionof{figure}{Kết quả mô hình không sử dụng metadata}
    \label{fig:model_khong_metadata}
\end{center}

\begin{center}
    \includegraphics[width=0.6\linewidth]{Hinhve/C5/mo_hinh_meta+tune.jpg}
    \captionof{figure}{Kết quả mô hình sử dụng metadata}
    \label{fig:model_meta+tune}
\end{center}

\textit{Giải thích:} 
Trong một số trường hợp, việc sử dụng metadata (user/item features) có thể khiến mô hình hoạt động kém hơn vì:

\begin{itemize}
    \item Dữ liệu metadata chưa đủ đặc trưng hoặc bị nhiễu: Ví dụ như đặc trưng của sản phẩm quá chung chung (như category phổ biến) hoặc user profile không phân biệt rõ nhóm người dùng.
    \item Mô hình bị overfitting do số lượng đặc trưng lớn nhưng phân bố không đồng đều, đặc biệt là khi số lượng user/item nhỏ so với feature dimension.
    \item Ma trận tương tác vốn đã rất đầy đủ và mạnh, nên việc thêm metadata đôi khi lại làm "loãng" tín hiệu gốc từ hành vi người dùng.
\end{itemize}
Do đó, dù về lý thuyết mô hình hybrid có thể cải thiện kết quả, nhưng thực tế tùy vào chất lượng feature, metadata không phải lúc nào cũng nâng cao hiệu suất.

 Đây là một phát hiện thực tiễn quan trọng, cho thấy không phải lúc nào thêm nhiều đặc trưng cũng cải thiện kết quả. Việc đánh giá hiệu quả các nguồn dữ liệu đầu vào là rất cần thiết trong các hệ thống gợi ý.
 Hình \ref{fig:suggestScreen} đã nêu ở chương 4 minh họa cho kết quả triển khai mô hình gợi ý này trên giao diện mobile của người dùng

\section{Triển khai tính năng "Chatbot tự động hóa quy trình mua sắm"}

\subsection{Đặt vấn đề}

Trong các hệ thống thương mại điện tử hiện nay, người dùng thường phải thực hiện nhiều thao tác thủ công lặp đi lặp lại để hoàn tất một quy trình mua sắm cơ bản: từ việc tìm kiếm sản phẩm, lọc lựa chọn phù hợp, thêm vào giỏ hàng, đến kiểm tra và xác nhận đặt hàng. Điều này khiến trải nghiệm mua sắm trở nên rườm rà, đặc biệt là trên thiết bị di động hoặc đối với người dùng không có nhiều thời gian hay không rành công nghệ.

Ngoài ra, một số người dùng có xu hướng cần trợ giúp trực tiếp hoặc thích tương tác bằng ngôn ngữ tự nhiên thay vì thao tác qua nhiều bước trên giao diện truyền thống. Trong khi đó, nhân sự hỗ trợ khách hàng lại không thể làm việc 24/7, dẫn đến tình trạng nghẽn tương tác, mất đơn hàng, hoặc bỏ lỡ nhu cầu tiềm năng.

Từ phía doanh nghiệp, việc không thể tự động hóa các quy trình cơ bản này làm tăng chi phí vận hành, chậm trễ trong việc phục vụ người dùng, và giảm hiệu quả trong việc chuyển đổi hành vi mua sắm. Đặc biệt, trong bối cảnh số lượng người dùng tăng nhanh, dữ liệu giao dịch ngày càng nhiều thì việc tối ưu hóa hành trình người dùng và giảm tải cho hệ thống là một yêu cầu cấp thiết.

Vì vậy, cần có một giải pháp thông minh, có khả năng giao tiếp bằng ngôn ngữ tự nhiên, hiểu được mục đích người dùng, và tự động thực hiện các thao tác như tìm kiếm, thêm sản phẩm, đặt hàng… chỉ thông qua hội thoại – giúp đơn giản hóa quy trình, cá nhân hóa trải nghiệm và nâng cao hiệu quả vận hành toàn hệ thống.

\subsection{Giải pháp}

Với những lợi ích rõ rệt như tiết kiệm thời gian thao tác, nâng cao trải nghiệm người dùng và tối ưu hiệu suất vận hành, em lựa chọn triển khai chatbot sử dụng nền tảng mã nguồn mở \textbf{Rasa} – một framework mạnh mẽ hỗ trợ xử lý ngôn ngữ tự nhiên (NLP) và xây dựng chatbot có khả năng thực hiện hành động cụ thể theo ngữ cảnh.

\textbf{Kiến trúc hệ thống chatbot gồm các thành phần chính sau:}
\begin{itemize}
    \item \textbf{NLU (Natural Language Understanding)}: Phân tích câu nhập của người dùng để xác định \textit{intent} (mục đích) và \textit{entities} (các thực thể liên quan).
    \item \textbf{Core}: Quản lý trạng thái hội thoại, dựa vào intent và lịch sử trò chuyện để quyết định hành động tiếp theo.
    \item \textbf{Actions Server}: Xử lý các hành động tùy chỉnh phức tạp, ví dụ như gọi API để tìm sản phẩm, thêm vào giỏ hàng hoặc hiển thị giỏ hàng, được viết trong tệp \texttt{actions.py}.
    \item \textbf{Front-end Flutter App}: Gửi câu thoại của người dùng lên server Rasa và hiển thị các phản hồi (bao gồm cả dạng button) từ chatbot.
\end{itemize}

\textbf{Cấu trúc thư mục chính của chatbot Rasa:}
\begin{itemize}
    \item \texttt{nlu.yml}: Chứa tập dữ liệu huấn luyện NLU, định nghĩa các intent như \texttt{greet}, \texttt{search\_product}, \texttt{add\_to\_cart}, \texttt{place\_order}, v.v.
    \item \texttt{domain.yml}: Định nghĩa toàn bộ intents, entities, actions, slots, responses (bao gồm cả trả về button), và các session rule.
    \item \texttt{stories.yml} và \texttt{rules.yml}: Mô tả luồng hội thoại bằng cách ghi lại các kịch bản mẫu và quy tắc hội thoại.
    \item \texttt{actions.py}: Chứa các hàm xử lý logic như gọi API sản phẩm, quản lý giỏ hàng, xử lý đặt hàng,... Các hàm này có thể trả về tin nhắn dạng JSON (custom payload) để hiện các quick reply buttons trên UI.
\end{itemize}

\textbf{Flow hoạt động tổng thể:}

\begin{itemize}

    \item Người dùng nhập một câu như: \texttt{"Mình muốn mua găng tay"}.
    \item Rasa NLU phân tích và gán intent là \texttt{search\_product}, đồng thời trích xuất entity là \texttt{găng tay}.
    \item Rasa Core xác định hành động tương ứng là \texttt{action\_search\_product}.
    \item Tại \texttt{actions.py}, action này gọi đến API tìm kiếm sản phẩm với từ khóa "găng tay", sau đó trả về danh sách kết quả dưới dạng json, sau đó sẽ trích xuất trường title, giơi hạn trả về 50 sản phẩm, đánh số thứ tự sản phẩm để cho bước tiếp theo như thêm vào giỏ hàng dễ dàng nhập "thêm 2 cái số 3 vào giỏ hàng" thay vì phải nhập cả tên sản phẩm rất dài. Ngoài ra thì response cũng bao gồm trường button với title và payload phục vụ cho việc hiển thị quick button ở trên giao diện, tối ưu trải nghiệm người dùng, giúp họ không cần nhập thủ công mà chỉ cần click button gợi ý trên giao diện là bot sẽ tự động thực thi hành động tương ứng, vô cùng nhanh và tiện lợi
    \item Trên giao diện Flutter, hệ thống gửi câu lệnh mà người dùng nhập vào lên rasa server, nhận response trả về, bóc tách các trường để xử lý logic và hiển thị text, hiển thị quick button lên giao diện 
    \item Tương tự, các flow như \texttt{hiển thị giỏ hàng}, \texttt{đặt hàng}, \texttt{hỏi tiếp sản phẩm khác}, \texttt{xin chào/tạm biệt}... đều được triển khai dưới dạng intent riêng với action tùy chỉnh tương ứng.
\end{itemize}
Sau khi đã khai báo và hoàn thiện các file trong rasa, chỉ cần gõ lệnh rasa run để build model, mất khoảng 2-3 phút, sau đó chạy lệnh 
rasa run actions để run actions server và rasa run để chạy rasa server, rasa server sẽ tự động gọi đến action-server qua port 5055. Sau khi đã build rasa server thành công, flutter mobile sẽ gọi đến server bằng https request để gửi tin nhắn và nhận phản hồi
\subsection{Kết quả}
Người dùng có giao diện chat bot đẹp mắt để trò chuyện với bot, tự động hóa quy trình mua sắm ngay trong chat page, có thể nhập hoặc đơn giản hơn là chỉ với vài click là có thể ra lệnh cho chatbot thực thi hành động ngay lập tức. Luồng trò chuyện hoàn chỉnh với chatbot 
được tôi demo qua các hình ảnh ở dưới

% ======= Hàng 1: 3 ảnh ========
\begin{minipage}[t]{0.3\textwidth}
    \centering
    \includegraphics[width=\linewidth]{Hinhve/C5/chat_bot/hello.jpg}
    \captionof{figure}{Người dùng bắt đầu cuộc hội thoại bằng lời chào, bot phản hồi ngay bằng các quick button gợi ý khám phá}
\end{minipage}
\hfill
\begin{minipage}[t]{0.3\textwidth}
    \centering
    \includegraphics[width=\linewidth]{Hinhve/C5/chat_bot/list_product.jpg}
    \captionof{figure}{Khi người dùng bấm khám phá 1 sản phẩm, bot hiển thị danh sách sản phẩm tìm thấy}
\end{minipage}
\hfill
\begin{minipage}[t]{0.3\textwidth}
    \centering
    \includegraphics[width=\linewidth]{Hinhve/C5/chat_bot/quick_button_goi_y_kham_pha.jpg}
    \captionof{figure}{Bot đưa ra các gợi ý là các quick button để người dùng có thể thêm nhanh sản phẩm vào giỏ hàng hoặc nhập câu lệnh tương tự theo mẫu hoặc xem giỏ hàng }
\end{minipage}

\vspace{0.5cm}

% ======= Hàng 2: 3 ảnh ========
\begin{minipage}[t]{0.3\textwidth}
    \centering
     \includegraphics[width=\linewidth]{Hinhve/C5/chat_bot/cart.jpg}
    \captionof{figure}{Bot điều hướng đến giao diện giỏ hàng khi người dùng chọn thêm vào giỏ hàng}
\end{minipage}
\hfill
\begin{minipage}[t]{0.3\textwidth}
    \centering
    \includegraphics[width=\linewidth]{Hinhve/C5/chat_bot/order.jpg}
    \captionof{figure}{Bot điều hướng đến xem đơn hàng và đặt hàng nếu người dùng chọn xem đơn hàng}
\end{minipage}
\hfill
\begin{minipage}[t]{0.3\textwidth}
    \centering
      \includegraphics[width=\linewidth]{Hinhve/C5/chat_bot/tim_them_san_pham.jpg}
    \captionof{figure}{Bot hiển thị danh sách các mục để người dùng có thể quick button click và xem thêm các sản phẩm khác tiếp theo}
\end{minipage}

\vspace{0.5cm}

% ======= Hàng 3: 3 ảnh ========
\begin{minipage}[t]{0.3\textwidth}
    \centering
   \includegraphics[width=\linewidth]{Hinhve/C5/chat_bot/tim_san_pham_thu_cong.jpg}
    \captionof{figure}{Ngoài việc bấm quickly button, người dùng vẫn có thể tự nhập để tìm kiếm sản phẩm mình muốn ngoài những gợi ý của bot}
\end{minipage}
\hfill
\begin{minipage}[t]{0.3\textwidth}
    \centering
    \includegraphics[width=\linewidth]{Hinhve/C5/chat_bot/cart_thu_cong.jpg}
    \captionof{figure}{Ngoài việc bấm quickly button, người dùng vẫn có thể tự nhập để thêm sản phẩm và số lượng mình thích vào giỏ hàng ngoài những gợi ý có sẵn của bot}
\end{minipage}
\hfill
\begin{minipage}[t]{0.3\textwidth}
    \centering
    \includegraphics[width=\linewidth]{Hinhve/C5/chat_bot/tam_biet.jpg}
    \captionof{figure}{Bot cảm ơn người dùng và kết thúc phiên trò chuyện.}
\end{minipage}

\section{Triển khai tính năng "Săn sale cùng bạn bè và chia sẻ nhanh qua mạng xã hội"}
\label{section:5.4}
\subsection{Đặt vấn đề}

Trong thời đại số hiện nay, hành vi tiêu dùng của người dùng đang dần thay đổi theo hướng cộng đồng và tương tác cao. Người dùng không còn đơn thuần mua sắm một cách cá nhân, mà có xu hướng tham khảo ý kiến bạn bè, chia sẻ khuyến mãi hấp dẫn và cùng nhau tìm kiếm các cơ hội mua sắm tiết kiệm.

Tuy nhiên, phần lớn các hệ thống hiện nay vẫn chưa thực sự tập trung phát triển tính năng mua sắm theo nhóm kết hợp chia sẻ qua mạng xã hội. Dù người dùng có nhu cầu rõ rệt trong việc kết nối bạn bè và cùng nhau săn sale, đa số nền tảng vẫn chỉ dừng lại ở các hình thức chia sẻ thủ công hoặc thiếu sự tích hợp trực tiếp với các kênh mạng xã hội quen thuộc. Điều này khiến người dùng thường bỏ lỡ các cơ hội mua sắm theo nhóm do không kịp mời bạn bè tham gia, hoặc gặp khó khăn trong việc tìm và kết nối với những người có nhu cầu tương tự trong khoảng thời gian giới hạn.


Đây chính là khoảng trống lớn mà ứng dụng của em muốn khai thác: một nền tảng nơi người dùng có thể dễ dàng lập nhóm, mời bạn bè và chia sẻ thông tin săn sale chỉ với vài thao tác đơn giản qua các kênh mạng xã hội quen thuộc như Messenger, Zalo, Discord,... Việc mua sắm không chỉ còn là hành động cá nhân, mà trở thành một trải nghiệm kết nối, tiết kiệm và đầy hứng khởi.

Nhu cầu chia sẻ, kết nối và mua sắm thông minh đang trở thành xu hướng tất yếu. Nếu ứng dụng không kịp thời bắt nhịp với xu thế này, rất có thể sẽ đánh mất nhóm người dùng trẻ năng động – những người luôn muốn vừa mua sắm vừa kết nối xã hội. Vì vậy, một tính năng săn sale nhóm kết hợp chia sẻ nhanh qua mạng xã hội là không chỉ cần thiết mà còn là một điểm khác biệt chiến lược, tạo lợi thế cạnh tranh cho sản phẩm.

\subsection{Giải pháp}
Đầu tiên, trong mục My Sale Groups của mục Profile trong ứng dụng, khi người dùng click vào sẽ xem danh sách các nhóm họ đã tạo hoặc đã tham gia theo lời mời của bạn bè. Nếu chưa có nhóm nào thì họ sẽ điều hướng đến màn hình tạo group để tạo group săn sale mới. Màn hình này sẽ có form để người dùng nhập vào các thông tin như tên nhóm, số lượng người tối thiểu của nhóm để kích hoạt voucher, và phạm vi voucher áp dụng: các sản phẩm thuộc cùng 1 shop hoặc 1 brand hoặc 1 category nhất định. Sau khi đã chọn đầy đủ thành công và click nút tạo nhóm, hệ thống sẽ tạo và lưu nhóm mới tạo lên database firestore, tự tính toán voucher với mức giảm giá dựa trên số người và phạm vi được lựa chọn. Sau khi tạo thành công hệ thống gửi thông báo đã tạo nhóm thành công đến cho người dùng, kèm các thông tin cơ bản của nhóm, và cập nhật lại danh sách nhóm. Với mỗi nhóm người dùng có thể xem thông tin chi tiết của nhóm đó, gồm tên, phạm vi, mức sale, số thành viên hiện tại và số thành viên cần đạt, thời gian tạo, thời gian hết hạn, trạng thái nhóm, thanh tiến độ và đồng hồ đếm ngược 24h. 

Để mời thêm bạn bè của mình tham gia nhóm để nhóm đạt đủ KPI thành viên trước deadline, em có xây dựng 2 cách để mời khác nhau:

\begin{itemize}
    \item Thứ nhất là mời với những bạn bè đã kết bạn với nhau trong hệ thống. Với cách mời này, người dùng sẽ chọn và được điều hướng đến màn hình mời bạn bè tham gia nhóm. Màn hình sẽ hiển thị danh sấch bạn bè hiện tại của người dùng, với mỗi người bạn sẽ có kèm một nút mời. Khi người dùng click vào nút này, sẽ gửi request đến bạn của mình, nếu request được gửi thành công sẽ có thông báo từ hệ thống gửi đến người dùng đã gửi thành công và đến người bạn đó là có 1 lời mời tham gia nhóm mới. Người bạn đó khi nhận được thông báo sẽ vào màn hình danh sách request để kiểm tra, Ở đây người đó có thể chấp nhận hoặc từ chối lời mời. Nếu chấp nhận thì người đó sẽ được hệ thống thêm vào nhóm, các thông tin về số thành viên hiện tại, tiến độ, trạng thái của nhóm cũng sẽ được tự động cập nhật tức thời. Nếu số lượng người tham gia đạt đến yêu cầu và vẫn trước deadline thì tất cả các thành viên trong nhóm sẽ đều nhận được thông báo nhóm đã tạo thành công, đã nhận được voucher săn sale của nhóm và tự động thêm vào mục các voucher người đó đã thu thập. Còn nếu người đó chọn từ chối thì cũng sẽ có thông báo gửi đến người dùng là lời mời của bạn đã bị bạn bè từ chối.
    \item Thứ hai là gửi link hoặc QR chia sẻ qua mạng xã hội. Với giải pháp này thì có thể mời bất kì ai, kể cả đã kết bạn hay chưa. Người dùng chọn mục chia sẻ qua mạng xã hội, sẽ có 1 pop up hiện lên với hai lựa chọn chia sẻ link liên kết hoặc chia sẻ QR. Sau đó sẽ có popup hiện lên cho người đó lựa chọn nền tảng muốn chia sẻ, chính là các nền tảng mạng xã hội mà người đó đã cài đặt như: messenger, zalo, facebook, instagram... Khi bạn bè nhận được link hoặc QR, bấm vào link hoặc quét QR, tùy thuộc vào việc người đó chưa tải và đăng nhập ứng dụng thì hệ thống sẽ điều hướng người đó đến màn hình cho phép nhập email của họ và sẽ gửi link tải ứng dụng qua email. Người đó có thể mở email, tải ứng dụng về máy, tạo tài khoản và đăng nhập vào hệ thống. Việc chia sẻ link tải và invite qua email này là nhờ em đã tích hợp dịch vụ \cite{FirebaseAppDistribution}. Trường hợp người đó đã tải và đã đăng nhập hệ thống rồi, thì khi mở link và quét QR, hệ thống tự điều hướng đến màn hình home và popup hiển thị thông tin nhóm, nút chọn tham gia nhóm. Khi người dùng chọn tham gia nhóm thì không cần chờ chủ nhóm phê duyệt mà có thể tham gia nhóm ngay lập tức.    
\end{itemize}

Để phục vụ cho tính năng mời bạn bè đã kết bạn trong hệ thống, em cũng đã phát triển \textbf{tính năng kết bạn trong nhóm}. Ở mục My Friends trong mục Profile, người dùng có thể xem danh sách bạn bè hiện tại của mình trong ứng dụng. Nếu chưa có hoặc muốn kết thêm bạn, người đó sẽ vào màn hình kết bạn, ở đây sẽ hiển thị danh sách các người dùng trong hệ thống, có nút tìm kiếm để tìm một người mà họ quen theo email hoặc username của họ. Nếu muốn kết bạn, người đó sẽ click vào nút add friend, hệ thống sẽ tự động gửi request đến người kia, và gửi thông báo là bạn có một lời mời kết bạn mới. Người nhận sẽ vào check danh sách lời mời kết bạn, cũng có thể chấp nhận hoặc từ chối. Nếu chấp nhận thì cả hai sẽ trở thành bạn bè, và gửi thông báo cho người gửi lời mời. Còn nếu người nhận từ chối thì cũng sẽ có thông báo đã bị từ chối và có thể gửi lời mời khác.

Khi các nhóm đã qua deadline mà chưa tập hợp đủ người, thì nhóm đó sẽ tự động chuyển sang trạng thái expired, và chủ nhóm sẽ không thể mời thêm người vào nhóm nữa.

\subsection{Kết quả}
Ở chương 4, em đã có minh họa một số kết quả của tính năng này, đó là màn hình thêm bạn bè \ref{fig:addFriendScreen}, màn hình chấp nhận tham gia nhóm \ref{fig:approveGroupRequestScreen}, màn hình tạo nhóm săn sale \ref{fig:createSaleGroup}, màn hình chi tiết nhóm săn sale \ref{fig:groupDetail}, màn hình mời bạn bè tham gia nhóm săn sale \ref{fig:inviteFriend}.

Ngoài các màn hình đó, em sẽ minh họa tính năng chia sẻ qua mạng xã hội:

% ======= Hàng 1: 3 ảnh ========
\begin{minipage}[t]{0.3\textwidth}
    \centering
    \includegraphics[width=\linewidth]{Hinhve/C5/san_sale/chia_sẻ.jpg}
    \captionof{figure}{Người dùng chọn chức năng chia sẻ qua mạng xã hội, có 2 lựa chọn là gửi link hoặc gửi QR}
\end{minipage}
\hfill
\begin{minipage}[t]{0.3\textwidth}
    \centering
    \includegraphics[width=\linewidth]{Hinhve/C5/san_sale/qr.jpg}
    \captionof{figure}{Khi người dùng chọn QR, hệ thống sẽ gen sẵn ảnh QR thông tin nhóm để người đó chia sẻ}
\end{minipage}
\hfill
\begin{minipage}[t]{0.3\textwidth}
    \centering
    \includegraphics[width=\linewidth]{Hinhve/C5/san_sale/gui_mess.jpg}
    \captionof{figure}{Hình ảnh cả link và cả QR đã được chia sẻ qua messenger}
\end{minipage}

\vspace{0.5cm}

% ======= Hàng 2: 3 ảnh ========
\begin{minipage}[t]{0.3\textwidth}
    \centering
     \includegraphics[width=\linewidth]{Hinhve/C5/san_sale/tai_ung_dung.jpg}
    \captionof{figure}{Nếu bấm vào link hoặc quét QR mà người đó chưa cài app, sẽ yêu cầu họ nhập email để invite qua email}
\end{minipage}
\hfill
\begin{minipage}[t]{0.3\textwidth}
    \centering
    \includegraphics[width=\linewidth]{Hinhve/C5/san_sale/popup_tg.jpg}
    \captionof{figure}{Nếu bấm vào link hoặc quét QR mà người đó đã cài app, sẽ có popup mời họ tham gia}
\end{minipage}
\hfill
\begin{minipage}[t]{0.3\textwidth}
    \centering
      \includegraphics[width=\linewidth]{Hinhve/C5/san_sale/thong_bao.jpg}
    \captionof{figure}{Khi tập hợp đủ người, sẽ có thông báo nhóm tập hợp đủ người và thông báo gửi tới các thành viên đã nhận voucher thành công}
\end{minipage}

\section{Triển khai tính năng "Làm nhiệm vụ và đổi thưởng"}
\subsection{Đặt vấn đề}

Trong bối cảnh thị trường thương mại điện tử ngày càng cạnh tranh, việc thu hút người dùng mới đã trở nên khó khăn hơn bao giờ hết — nhưng giữ chân người dùng cũ và tăng mức độ tương tác với nền tảng còn là một thách thức lớn hơn. Người dùng ngày nay không chỉ kỳ vọng vào trải nghiệm mua sắm tiện lợi, mà còn mong muốn được tương tác, tham gia và nhận lại giá trị khi họ chủ động sử dụng ứng dụng thường xuyên.

Tuy nhiên, thực tế cho thấy phần lớn người dùng chỉ quay lại ứng dụng khi có nhu cầu mua sắm cụ thể, sau đó rời đi mà không duy trì sự hiện diện thường xuyên. Điều này dẫn đến tỷ lệ người dùng hoạt động hàng ngày (DAU) thấp, thời gian phiên sử dụng ngắn và khả năng tạo ra doanh thu bị hạn chế. Các nền tảng thiếu cơ chế duy trì thói quen sử dụng, khuyến khích hành vi tích cực, hoặc tạo cảm giác “được tưởng thưởng” cho những người dùng trung thành.

Trong xu hướng game hoá (gamification) ngày càng được áp dụng rộng rãi, việc tích hợp các yếu tố như nhiệm vụ hàng ngày, thử thách theo thời gian, hệ thống điểm thưởng và quy đổi quà tặng không chỉ giúp tăng tính hấp dẫn mà còn tạo động lực để người dùng quay lại và tương tác thường xuyên. Đây không chỉ là một cách để làm cho trải nghiệm mua sắm trở nên thú vị hơn, mà còn là chiến lược hiệu quả để tăng mức độ trung thành, kéo dài vòng đời người dùng và thúc đẩy doanh thu.

Chính vì vậy, việc triển khai một hệ thống nhiệm vụ kết hợp đổi thưởng là một giải pháp cấp thiết để nâng cao sự gắn bó của người dùng, gia tăng tần suất sử dụng ứng dụng, đồng thời tạo ra một lợi thế cạnh tranh bền vững trên thị trường.

\subsection{Giải pháp}
Từ vấn đề trên, em đã xây dựng tính năng làm nhiệm vụ tích điểm thưởng và tính năng đổi điểm lấy voucher.

Về các nhiệm vụ, thử thách cho người dùng, thì đầu tiên là nhiệm vụ điểm danh hàng ngày. Em sẽ tạo chuỗi thử thách 7 ngày điểm danh liên tục, chuỗi điểm danh liên tiếp càng dài thì ngày tiếp theo sẽ được càng nhiều điểm hơn. Nếu người đó bị miss điểm danh ngày nào thì chuỗi tự reset và quay về từ ngày 1. Nếu đã hoàn thành điểm danh liên tục cả 7 ngày cũng sẽ tự reset để bắt đầu một chuỗi điểm danh mới.

Nhóm nhiệm vụ thứ hai là thử thách thực sự với người dùng, khi phải thực hiện đủ số lượng nhiệm vụ yêu cầu trong một thời gian quy định. Người dùng vào phần My Missions của mục Profile trong app, sẽ thấy các nhiệm vụ được hệ thống tạo sẵn, và trạng thái nhiệm vụ nào mình chưa làm, nhiệm vụ nào đang làm, đã làm hoặc đã hết hạn. Khi xem chi tiết một nhiệm vụ sẽ có ảnh minh họa nhiệm vụ, tên, mô tả chi tiết nhiệm vụ, điểm thường, deadline, và đồng hồ đếm ngược thời gian còn lại của nhiệm vụ, cũng như thanh tiến trình hiển thị nhiệm vụ đã hoàn thành được bao nhiêu \% . Khi người dùng bắt đầu một nhiệm vụ, hệ thống sẽ tự tracker toàn ứng dụng để theo dõi khi nào người dùng thực hiện hành động nào, nếu hành động đó trùng với hành động mà nhiệm vụ đang được tính thời gian thì nó sẽ ghi nhận và tự động cập nhật tiến độ của nhiệm vụ. Để theo dõi tiến độ hoàn thành nhiệm vụ một cách tự động và chính xác, hệ thống được thiết kế theo hướng linh hoạt và mở rộng dễ dàng. Khi người dùng bắt đầu một nhiệm vụ, hệ thống sẽ tự động theo dõi toàn bộ ứng dụng để nhận diện khi người dùng thực hiện các hành động có liên quan. Nếu hành động đó trùng với loại nhiệm vụ đang được kích hoạt và tính thời gian, hệ thống sẽ tự động ghi nhận và cập nhật tiến độ.
\begin{center}
    \includegraphics[width=0.9\linewidth]{Hinhve/C5/Class_diagram_mission_tracking.jpg}
    \captionof{figure}{Biểu đồ class cho tính năng Mission Tracking}
    \label{fig:class_dg_mission_tracking}
\end{center}
Hình \ref{fig:class_dg_mission_tracking}
 minh họa biểu đồ lớp của hệ thống theo dõi nhiệm vụ. Các lớp chính trong hệ thống được thiết kế rõ ràng, tuân thủ nguyên tắc thiết kế hướng đối tượng và đảm bảo khả năng mở rộng linh hoạt:

\begin{itemize}
    \item \textbf{MissionTracker}:
    \begin{itemize}
        \item Là một \textbf{service trung tâm} của hệ thống, đóng vai trò như bộ điều phối logic xử lý nhiệm vụ khi người dùng thực hiện hành động tương ứng.
        \item Được triển khai theo \textbf{Singleton pattern}, đảm bảo chỉ có một thể hiện duy nhất trong toàn bộ ứng dụng. Điều này giúp dễ dàng gọi hàm track ở bất kỳ đâu trong ứng dụng cần tracking như ở màn hình xem sản phẩm, xem thương hiệu, hay khi người dùng đặt đơn hàng.
        \item Chứa một \textbf{Map các handler} – mỗi handler xử lý một loại nhiệm vụ nhất định, ánh xạ theo enum `MissionType`. Điều này giúp hệ thống xử lý linh hoạt đa dạng nhiệm vụ mà không cần thay đổi cấu trúc lõi.
        \item Phương thức track truy vấn Firestore để tìm các nhiệm vụ đang hoạt động thuộc type tương ứng, sau đó gọi đến handler tương ứng để xử lý logic cập nhật nhiệm vụ.
    \end{itemize}

    \item \textbf{MissionHandler (abstract class)}:
    \begin{itemize}
        \item Là lớp trừu tượng, định nghĩa phương thức `handle(...)` như một hợp đồng chung mà tất cả các lớp con phải hiện thực.
      
        \item Phụ thuộc UserMissionModel – đây là model đại diện cho một nhiệm vụ cụ thể mà người dùng đang tham gia. Handler sử dụng model này để lấy thông tin chi tiết và cập nhật tiến độ nhiệm vụ.
    \end{itemize}

    \item \textbf{ExploreMissionHandler, OrderMissionHandler}:Là các lớp cụ thể kế thừa từ MissionHandler, hiện thực phương thức handle(...) tương ứng với từng loại hành động: xem nội dung, đặt hàng, v.v.
\end{itemize}

\noindent
Hệ thống trên sở hữu nhiều điểm mạnh nổi bật về thiết kế:

\begin{itemize}
    
    \item \textbf{Tính mở rộng cao và tuân thủ nguyên lý \textit{Open-Closed}}:
    \begin{itemize}
        \item Khi cần thêm một loại nhiệm vụ mới (ví dụ: "chia sẻ sản phẩm lên mạng xã hội"), em chỉ cần:
        \begin{enumerate}
            \item Định nghĩa một giá trị mới trong enum \texttt{MissionType}.
            \item Tạo một lớp mới kế thừa từ \texttt{MissionHandler}, ví dụ \texttt{ShareMissionHandler}.
            \item Đăng ký handler mới vào map \_handlers trong \texttt{MissionTracker}.
        \end{enumerate}
        \item Toàn bộ quy trình xử lý tracking vẫn hoạt động như cũ, không cần chỉnh sửa lại các class hiện tại. Đây là minh chứng rõ ràng cho việc hệ thống \textbf{mở với mở rộng, đóng với chỉnh sửa}.
    \end{itemize}
    
    \item \textbf{Áp dụng mô hình Strategy Pattern} thông qua việc phân phối nhiệm vụ đến handler phù hợp với chiến lược xử lý riêng biệt.
    
    \item \textbf{Tự động hóa toàn diện}: Người dùng không cần thao tác thủ công để cập nhật tiến độ nhiệm vụ. Mỗi hành động được hệ thống tự động ghi nhận và xử lý đúng với logic tương ứng.
    
    \item \textbf{Tuân thủ nguyên lý Single Responsibility (SRP)}: \texttt{MissionTracker} chỉ điều phối nhiệm vụ, còn các \texttt{MissionHandler} xử lý logic nhiệm vụ cụ thể, đảm bảo sự rõ ràng và dễ bảo trì.
\end{itemize}

Nhờ cách tổ chức module hợp lý, hệ thống có thể dễ dàng mở rộng để hỗ trợ hàng chục, hàng trăm loại nhiệm vụ khác nhau trong tương lai, mà vẫn đảm bảo hiệu năng, tính ổn định và khả năng kiểm soát logic một cách rõ ràng.

Khi người dùng hoàn thành một nhiệm vụ, thì sẽ có thông báo nhiệm vụ đã hoàn thành, kèm số điểm nhận được, và điểm thưởng sẽ tự động được tăng lên. 

Hệ thống cũng có tích hợp tính năng đổi thưởng, khi hệ thống cũng sẽ quy định sẵn một số loại voucher có thể quy đổi bằng điểm thưởng, gồm thông tin voucher, mức điểm thưởng yêu cầu và nút quy đổi. Khi người dùng click nút quy đổi thì tự động voucher đó sẽ được nhận và điểm thưởng của người dùng sẽ bị trừ tương ứng.

Ngoài ra, ở phần web admin quản trị viên có thể quản lý CRUD các nhiệm vụ trong hệ thống, cũng như xem thống kê các nhiệm vụ theo type, theo lịch sử người dùng đã thực hiệnhiện

\subsection{Kết quả}
Ở chương 4, em đã minh họa các tính năng điểm danh hàng ngày \ref{fig:dailyCheckin} và màn hình nhiệm vụ chi tiết \ref{fig:missionDetail}. Ở chương này, em sẽ minh họa tiếp giao diện của các màn hình khác liên quan đến nhóm chức năng này

\begin{minipage}[t]{0.3\textwidth}
    \centering
    \includegraphics[width=\linewidth]{Hinhve/C5/danh_sach_nhiem_vu.jpg}
    \captionof{figure}{Màn hình danh sách các nhiệm vụ có sẵn trong hệ thống và trạng thái của người dùng đã hay chưa xử lý}
\end{minipage}
\hfill
\begin{minipage}[t]{0.3\textwidth}
    \centering
    \includegraphics[width=\linewidth]{Hinhve/C5/doi_thuong.jpg}
    \captionof{figure}{Màn hình quy đổi điểm thưởng lấy voucher}
\end{minipage}
\hfill
\begin{minipage}[t]{0.3\textwidth}
    \centering
    \includegraphics[width=\linewidth]{Hinhve/C5/thong_bao.jpg}
    \captionof{figure}{Thông báo khi người dùng hoàn thành 1 nhiệm vụ và nhận được điểm thưởng tương ứng}
\end{minipage}

\vspace{0.5cm}
\begin{center}
    \includegraphics[width=0.9\linewidth]{Hinhve/C5/CRUD_MISSSION.jpg}
    \captionof{figure}{Màn hình web admin CRUD nhiệm vụ}
    \label{fig:crud_mission}
\end{center}
Hình \ref{fig:crud_mission} hiển thị màn hình quản trị viên CRUD nhiệm  vụ, ngoài ra có thể search kết hơp lọc nhiệm vụ theo type

\begin{center}
    \includegraphics[width=0.9\linewidth]{Hinhve/C5/thong_ke_type.jpg}
    \captionof{figure}{Màn hình thống kê theo type}
    \label{fig:thong_ke_nhiem_vu_theo_type}
\end{center}

\begin{center}
    \includegraphics[width=0.9\linewidth]{Hinhve/C5/thong_ke_by_use.jpg}
    \captionof{figure}{Màn hình thống kê theo lịch sử thực hiện nhiệm vụ của người dùng}
    \label{fig:thong_ke_nhiem_vu_theo_use}
\end{center}

\section{Triển khai hệ thống "voucher, áp dụng và quản lý voucher"}
\subsection{Đặt vấn đề}

Trong bối cảnh thị trường thương mại điện tử ngày càng cạnh tranh, việc tạo ra các hình thức khuyến mãi thông minh và cá nhân hóa là một trong những yếu tố then chốt để thu hút và giữ chân người dùng. Một trong những công cụ quan trọng nhất hỗ trợ cho chiến lược này chính là hệ thống voucher – đóng vai trò như một cầu nối giữa ưu đãi từ nhà bán và hành vi mua sắm của người dùng.

Tuy nhiên, ở nhiều nền tảng hiện nay, hệ thống voucher vẫn còn tồn tại nhiều hạn chế. Việc tìm kiếm và áp dụng mã giảm giá thường yêu cầu người dùng phải nhớ mã, tự điền thủ công, hoặc phụ thuộc vào thông báo rời rạc từ hệ thống. Điều này gây bất tiện, đặc biệt trong các tình huống cần thao tác nhanh như trong giai đoạn thanh toán. Ngoài ra, người dùng cũng gặp khó khăn trong việc theo dõi các voucher đã thu thập, thời hạn sử dụng, và điều kiện áp dụng cụ thể cho từng đơn hàng. Tình trạng "quên dùng" voucher vẫn diễn ra thường xuyên, làm giảm hiệu quả của chương trình khuyến mãi, đồng thời ảnh hưởng đến trải nghiệm người dùng.

Về phía nhà quản trị hệ thống, việc không có một nền tảng quản lý voucher tập trung và trực quan khiến cho quá trình tạo mới, kiểm soát, phân tích hiệu quả, và đo lường hành vi người dùng trở nên thủ công, thiếu chính xác và khó mở rộng. Điều này làm giảm khả năng triển khai các chiến dịch khuyến mãi linh hoạt và mang tính cá nhân hóa theo từng đối tượng mục tiêu.

Chính vì vậy, nhu cầu về một hệ thống voucher toàn diện – từ việc thu thập, theo dõi, áp dụng cho đến quản lý và thống kê – là vô cùng cần thiết. Việc thiết kế một hệ thống như vậy không chỉ giúp tối ưu trải nghiệm người dùng khi mua sắm, mà còn mang lại công cụ mạnh mẽ để các nhà phát triển và quản trị viên dễ dàng điều phối chiến lược khuyến mãi một cách linh hoạt, minh bạch và hiệu quả.

\subsection{Giải pháp}
Để giải quyết vấn đề trên, đầu tiên trong ứng dụng của em khi người dùng mở mục My Vouchers trong phần Profile, người dùng sẽ xem được danh sách các voucher mà người đó chưa thu thập và vẫn còn hiệu lực trong hệ thống. Màn hình có giao diện thân thiện, đẹp mắt, có nút seach kết hợp với lọc theo từng loại voucher giúp người dùng nhanh chóng tìm thấy voucher mong muốn. Để tối ưu thời gian phản hồi, hệ thống cũng không hiển thị toàn bộ voucher mà áp dụng chiến lược phân trang, mỗi lần chỉ hiển thị 20 voucher đầu tiên, người dùng muốn xem thêm sẽ click nút loadmore sẽ lại load thêm 20 voucher tiếp theo. Với mỗi voucher sẽ có nút claim để người đó có thể thu thập và chuẩn bị sử dụng.

Thứ hai, từ màn hình danh sách voucher có thể điều hướng sang màn hình lịch sử voucher của người dùng. Màn hình này có giao diện gần tương tự màn hình danh sách ở trên, nhưng có chia làm 3 tab là voucher đã thu thập, voucher đã sử dụng và voucher đã hết hạn để người dùng dễ dàng quản lý việc sử dụng của mình

Thứ ba, ở màn hình checkout, trước khi click thanh toán có nút apply voucher, khi người dùng ấn nút này, thì sẽ hiển thị popup các loại voucher người đó đã thu thập, chưa sử dụng, còn thời hạn hiệu lực và đủ điều kiện áp dụng với điều kiện đơn hàng hiện tại. Khi người dùng chọn một voucher cụ thể và click button apply, hệ thống sẽ tự động tính toán được trừ của đơn hàng sau khi áp dụng, và hiển thị thông tin của voucher vừa áp dụng lên trên giao diện màn hình checkout ngay lập tức. Hệ thống cũng tự động kiểm tra nếu giá trị đơn hàng đã về 0 hoặc đã áp dụng 1 voucher 2 lần( ví dụ voucher freeship) cũng sẽ có cảnh báo và không tiến hành áp dụng để không lãng phí voucher

Thứ tư, ở phần web admin, quản trị viên có thể quản lý CRUD các voucher trong hệ thống, cũng như xem thống kê các voucher theo type, theo lịch sử người dùng đã áp dụng.

\subsection{Kết quả}
 Dưới đây em sẽ minh họa các giao diện liên quan đến nhóm chức năng này: màn hình lịch sử voucher, popup apply voucher và màn hình chekout khi áp dụng thành công 1 voucher

\begin{minipage}[t]{0.3\textwidth}
    \centering
    \includegraphics[width=\linewidth]{Hinhve/C5/voucher_history.jpg}
    \captionof{figure}{Màn hình danh lịch sử các voucher của người dùng có phân theo 3 nhóm khác nhau}
\end{minipage}
\hfill
\begin{minipage}[t]{0.3\textwidth}
    \centering
    \includegraphics[width=\linewidth]{Hinhve/C5/popup_voucher.jpg}
    \captionof{figure}{Pop up gợi ý những voucher đủ điều kiện để người dùng sử dụng với điều kiện đơn hàng hiện tại}
\end{minipage}
\hfill
\begin{minipage}[t]{0.3\textwidth}
    \centering
    \includegraphics[width=\linewidth]{Hinhve/C5/apply_voucher.jpg}
    \captionof{figure}{Màn hình checkout khi apply thành công 1 voucher}
\end{minipage}

\vspace{0.5cm}
\begin{center}
    \includegraphics[width=0.9\linewidth]{Hinhve/C5/CRUD_VOUCHER.jpg}
    \captionof{figure}{Màn hình web admin CRUD voucher}
    \label{fig:crud_voucher}
\end{center}
Hình \ref{fig:crud_voucher} hiển thị màn hình quản trị viên CRUD voucher, ngoài ra có thể search kết hơp lọc voucher theo type

\begin{center}
    \includegraphics[width=0.9\linewidth]{Hinhve/C5/voucher_by_type.jpg}
    \captionof{figure}{Màn hình thống kê voucher theo từng type}
    \label{fig:thong_ke_voucher_theo_type}
\end{center}

\begin{center}
    \includegraphics[width=0.9\linewidth]{Hinhve/C5/voucher_by_usage.jpg}
    \captionof{figure}{Màn hình thống kê voucher theo lịch sử áp dụng của người dùng}
    \label{fig:thong_ke_voucher_theo_usage}
\end{center}

Như vậy là ở chương 5 này, em đã trình bày những giải pháp và đóng góp nổi bật của mình trong ứng dụng. Sau đây, em xin phép được trình bày chương 6 - Kết luận và hướng phát triển.
\end{document}